Identifying mothers with substance use disorders during pregnancy may create an opportunity to provide medical care and social support that can improve maternal and infant health. But drug testing in medical settings is complicated. Maternal drug use is socially stigmatized, which may make discussions of drug testing uncomfortable and could make it difficult for health care providers to make objective decisions and recommendations. In some cases, drug testing may put the patient's medical best interests at odds with their best interests from a more holistic point of view, including legal or family stability implications. It is possible that some patients may refuse testing if they fear the test will be used against them. If they perceive medical staff to be untrustworthy, patients may cut back on medical care across the board. Furthermore, meconium and umbilical cord drug tests could be used to subject families to drug testing without their consent, raising further ethical dilemmas. These challenges are amplified by concerns that medical personnel use non-clinical factors like the patient's race to decide which patients to test for drug use during pregnancy. 

This study examined drug testing decisions in a large study population of $140,562$ newborns in Indiana. There was a large racial testing gap of 9.2 percentage points between Black and White infants, making the testing rate more than twice as high among Black infants compared with White infants. However, the positive testing rate among infants who were tested was about 20\% for both Black and White infants. Despite the gap in testing rates, health care providers appear to allocate tests so that the average risk of substance use among those tested is the same across racial groups. This does not rule out the possibility that providers engage in discriminatory testing at the margin. Depending on the distribution of risks in the patient population, it is possible that physicians test "deeper" into the low risk end of the distribution for one racial group even if the overall average risk among the tested is the same. 

Guidelines related to infant drug testing suggest that physicians should make testing decisions based on risks they observe from the health of the infant and the mother, including the medical history of the mother. In this paper, we used detailed EMR records to measure the main clinical risk factors that are supposed to guide physician testing decisions. Then we used an inverse propensity score weighting estimator to decompose the racial testing gap into a share that is accounted for by clinical risk factors and an unexplained component. We found that observed risk factors account for about $39\%$ of the Black-White testing gap. In principle, the remaining $61\%$ of the testing gap may reflect either risk factors that are observed by health care workers but are not observed in the EMR data, or differences in the way that health care workers interpret the same characteristics differently for White and Black samples. The fact that positive testing rates are about 20\% for both White and Black infants may suggest that medical staff observe some additional risk factors that are predictive of substance use. Importantly, observing differences in risk profiles between Black and White racial groups should not be interpreted as evidence that racial identity itself carries risk \citep{chokshi2022act}. Differences in risk exposure could be rooted in upstream disparities in housing, employment, victimization, and education that contribute to poorer health behaviors and outcomes among historically marginalized groups \citep{yearby2022structural}. 

Our analysis focused on comparisons of realized testing status (tested or not) and on the test outcomes among the tested. We did not study how patients were actually presented with the option to be tested or not tested, whether tested mothers explicitly consented to testing and knew that their infant was being tested, or whether untested mothers actually refused testing. These are important avenues for future research. It is possible -- for example -- that there are racial disparities in test refusal rates and in the way that a test refusal is respected or interpreted by medical personnel. We are not able to distinguish mothers who were presented with an opportunity to refuse or consent to infant testing, or to examine the post-decision health care utilization or outcomes of others who refused or consented to testing. Research on testing and consent procedures could help clarify whether and when drug testing undermines trust, and when it facilitates interventions that support the infant and mother. 

Our paper has several limitations. First, we cannot be sure that the electronic medical records contain all of the clinically relevant information underlying the test. It is possible that the physicians making testing decisions have additional information that justifies their decision and would explain the unexplained $61\%$ of the gap in testing rates. Second, although we are able to compare the positive testing rate for the overall sample of Black and White infants,  we are not able to measure positive testing rates for \textit{marginally tested patients} in the sense used in the literature on \textit{outcomes tests} for non-statistical discrimination \citep{ayres2002outcome, knowles2001racial, arnold2018racial}. The challenge is that in our data, there is no quasi-experimental instrumental variable that shifts a patient's probability of being tested independently of their risk of testing positive. Research that compares positive testing rates at the margin may shed light on the possible role of racial animus or false stereotypes in driving testing behavior. We hope to pursue these questions in future work using quasi-random variation in hospital staffing schedules to identify the marginally tested patients to better estimate the role of racial discrimination in rates of testing.

Despite these limitations, this study advances knowledge in several ways. We provide new evidence on clinician decisions on which infants to test for in-utero drug exposure using a large sample that covers a wide range of hospitals across Indiana. We document a large gap in testing rates and nearly identical positive testing rates. We also quantify the share of racial testing disparities that can be explained by differences in clinical risk factors. Our finding that significant disparities remain even after controlling for observable clinical risk factors calls for future investigations to better understand the process through which infants are tested, what clinicians do with the information they receive from a test, and how mothers respond to and are affected by the test and corresponding outcomes.


